\documentclass[10pt]{article}
\usepackage{geometry}
\geometry{a4paper, left=1.6cm, right=1.6cm, top=2.2cm, bottom=2cm, headheight=14pt}
\usepackage{fontspec}
\usepackage[spanish,es-noshorthands]{babel}
\usepackage{fancyvrb}
\usepackage[most]{tcolorbox}
\usepackage{tabularx}
\usepackage{enumitem}
\usepackage{xcolor}
\usepackage{parskip}
\usepackage[hidelinks]{hyperref}
\usepackage{fancyhdr}
\usepackage{titlesec}
\usepackage{multicol}

\setlength{\parindent}{0pt}
\setlist[enumerate]{leftmargin=*, itemsep=0.5mm, topsep=1mm}
\setlist[itemize]{leftmargin=*, itemsep=0.5mm, topsep=1mm}
\setlength{\columnsep}{0.7cm}
\setlength{\columnseprule}{0.3pt}

% =========================================
% COLORES
% =========================================
\definecolor{azuloscuro}{RGB}{0,51,102}
\definecolor{cajagris}{RGB}{244,244,244}
\definecolor{marcogris}{RGB}{170,170,170}

% =========================================
% FORMATO DE SECCIONES (regla en \linewidth, no \textwidth, por multicol)
% =========================================
\titleformat{\section}{\large\bfseries\color{azuloscuro}}{\thesection.}{0.4em}{}[\vspace{-0.35em}{\color{marcogris}\rule{\linewidth}{0.6pt}}]
\titlespacing*{\section}{0pt}{2.5mm}{1mm}
\titleformat{\subsubsection}[runin]{\small\bfseries\color{azuloscuro}}{}{0em}{}[.\quad]

% =========================================
% CAJAS
% =========================================
\newtcolorbox{cajacodigo}{
  colback=cajagris, colframe=marcogris, boxrule=0.4pt,
  arc=1mm, left=2mm, right=2mm, top=1mm, bottom=1mm, breakable
}
\newtcolorbox{formulabox}[1][]{
  colback=cajagris, colframe=azuloscuro, fonttitle=\bfseries\small,
  title=#1, breakable, arc=1mm, left=2mm, right=2mm, top=1mm, bottom=1mm
}

% =========================================
% ENCABEZADO
% =========================================
\pagestyle{fancy}
\fancyhf{}
\fancyfoot[C]{\small\thepage}
\renewcommand{\headrulewidth}{0.4pt}
\renewcommand{\footrulewidth}{0pt}
\lhead{\footnotesize Programación Python --- 4to Medio}
\rhead{\footnotesize Resumen: Fundamentos a Funciones}

\small

\begin{document}

\begin{center}
\Large\textbf{Resumen Python --- Fundamentos a Funciones}\\[0.5mm]
\footnotesize Datos y variables $\cdot$ Operadores $\cdot$ Entrada/Salida $\cdot$ Condicionales $\cdot$ Ciclos $\cdot$ Funciones
\end{center}

\vspace{2mm}

\begin{multicols}{2}

% =========================================
% 1. DATOS Y VARIABLES
% =========================================
\section*{1. Datos y variables}

Python guarda información en \textbf{variables}: un nombre que apunta a un valor de un tipo determinado.

\begin{itemize}
  \item \texttt{int} --- números enteros (\texttt{15}).
  \item \texttt{float} --- números decimales, con \textbf{punto}, nunca coma (\texttt{5.8}).
  \item \texttt{str} --- texto, entre comillas (\texttt{"Ana"}).
  \item \texttt{bool} --- \texttt{True} o \texttt{False}.
\end{itemize}

Asignación: \texttt{nombre\_variable = valor}. Nombres en \textbf{snake\_case}, significativos, sin palabras reservadas.

Conversión entre tipos: \texttt{int(x)}, \texttt{float(x)}, \texttt{str(x)}.

\begin{cajacodigo}
\begin{Verbatim}[fontsize=\footnotesize]
altura_persona = 175       # int
promedio_notas = 5.8       # float
nombre_estudiante = "Ana"  # str
esta_inscrito = True       # bool
\end{Verbatim}
\end{cajacodigo}

\begin{formulabox}[Palabras reservadas --- no usar como nombre de variable]
\footnotesize
\texttt{False, None, True, and, or, not, if, elif, else, for, while, def, return, import, break, continue, in}
\end{formulabox}

% =========================================
% 2. OPERADORES
% =========================================
\section*{2. Operadores}

\textbf{Aritméticos:} \texttt{+ \ - \ * \ / \ // \ \% \ **}\\
\texttt{//} es división entera (descarta el resto); \texttt{\%} es el resto (módulo); \texttt{**} es potencia.

\begin{cajacodigo}
\begin{Verbatim}[fontsize=\footnotesize]
17 // 5   >> 3
17 % 5    >> 2
2 ** 5    >> 32
\end{Verbatim}
\end{cajacodigo}

\textbf{De comparación} (siempre devuelven \texttt{bool}): \texttt{== \ != \ > \ < \ >= \ <=}

\textbf{Lógicos:}
\begin{itemize}
  \item \texttt{and} --- exige que \textbf{ambas} condiciones sean verdaderas.
  \item \texttt{or} --- basta con que \textbf{una} sea verdadera.
  \item \texttt{not} --- invierte el valor lógico.
\end{itemize}

\begin{cajacodigo}
\begin{Verbatim}[fontsize=\footnotesize]
(edad >= 12) and (edad <= 17)
\end{Verbatim}
\end{cajacodigo}

% =========================================
% 3. ENTRADA Y SALIDA
% =========================================
\section*{3. Entrada y salida}

\texttt{print(valor1, valor2, ...)} imprime en pantalla, separando con \textbf{comas} --- nunca concatenar con \texttt{+} y \texttt{str()}. \texttt{sep=} cambia el separador entre valores; \texttt{end=} cambia qué se imprime al final (por defecto, salto de línea).

\texttt{input("mensaje")} muestra el mensaje y devuelve \textbf{siempre texto} (\texttt{str}), aunque la persona escriba números. Para calcular con ese valor hay que convertirlo con \texttt{int()} o \texttt{float()}.

\begin{cajacodigo}
\begin{Verbatim}[fontsize=\footnotesize]
edad = int(input("¿Cuántos años tienes? "))
print("Tienes", edad, "años")
>> Tienes 15 años
\end{Verbatim}
\end{cajacodigo}

\textit{Convención de este resumen: el output de \texttt{print()} se muestra con \texttt{>>} en la línea siguiente.}

% =========================================
% 4. CONDICIONALES
% =========================================
\section*{4. Condicionales}

\texttt{if condición:} ejecuta el bloque indentado solo si la condición es \texttt{True}. \texttt{else:} ejecuta si no se cumplió. La \textbf{indentación} (4 espacios) define qué instrucciones pertenecen a cada bloque --- no es estética.

\texttt{elif condición:} agrega alternativas intermedias sin anidar --- equivale a "si no, y si se cumple esto otro".

\begin{cajacodigo}
\begin{Verbatim}[fontsize=\footnotesize]
if nota < 4.0:
    print("Insuficiente")
elif nota < 5.5:
    print("Regular")
else:
    print("Buena")
\end{Verbatim}
\end{cajacodigo}

\textbf{Anidados:} un \texttt{if}/\texttt{else} dentro de otro. Cada nivel agrega su propia indentación.

\begin{cajacodigo}
\begin{Verbatim}[fontsize=\footnotesize]
if es_socia:
    if compra >= 10000:
        print("Doble descuento")
    else:
        print("Descuento socia")
else:
    print("Sin descuento")
\end{Verbatim}
\end{cajacodigo}

% =========================================
% 5. CICLOS
% =========================================
\section*{5. Ciclos}

\texttt{for variable in range(...):} repite el bloque una vez por cada valor que entrega \texttt{range}. \texttt{range(inicio, fin, salto)} --- el \textbf{fin nunca se incluye}. Con un argumento, empieza en 0; con dos, el salto vale 1.

\begin{cajacodigo}
\begin{Verbatim}[fontsize=\footnotesize]
for numero in range(2, 11, 2):
    print(numero)
>> 2
>> 4
>> 6
>> 8
>> 10
\end{Verbatim}
\end{cajacodigo}

\textbf{\texttt{for} anidado:} un \texttt{for} dentro de otro (tablas, patrones) --- el interior completa todo su recorrido por cada paso del exterior.

\texttt{continue} salta el resto de esa iteración y sigue con la siguiente. \texttt{break} termina el ciclo completo de inmediato.

\begin{cajacodigo}
\begin{Verbatim}[fontsize=\footnotesize]
for numero in range(1, 8):
    if numero == 5:
        break
    if numero % 2 == 0:
        continue
    print(numero)
>> 1
>> 3
\end{Verbatim}
\end{cajacodigo}

\texttt{while condición:} repite \textbf{mientras} la condición sea \texttt{True} --- se usa cuando no se sabe de antemano cuántas repeticiones habrá. Hay que actualizar la variable de control dentro del bloque, o el ciclo nunca termina.

\begin{cajacodigo}
\begin{Verbatim}[fontsize=\footnotesize]
suma = 0
while suma < 100:
    suma = suma + 10
print("Suma final:", suma)
>> Suma final: 100
\end{Verbatim}
\end{cajacodigo}

\texttt{while True:} junto a un \texttt{if} + \texttt{break} interno crea un ciclo que se repite hasta que se cumple una condición decidida dentro del propio bloque.

% =========================================
% 6. FUNCIONES
% =========================================
\section*{6. Funciones}

\texttt{def nombre\_funcion(parametro1, parametro2):} define la función; el cuerpo va indentado debajo. \textbf{Definir la función todavía no la ejecuta} --- solo queda "guardada la receta" hasta que alguien la llama.

\begin{cajacodigo}
\begin{Verbatim}[fontsize=\footnotesize]
def calcula_total(precio, cantidad):
    return precio * cantidad
\end{Verbatim}
\end{cajacodigo}

\textbf{Parámetro:} variable que la función espera recibir (en la definición). \textbf{Argumento:} el valor real que se entrega al llamarla.

\texttt{return valor} entrega el resultado a quien llamó la función y \textbf{corta la ejecución ahí mismo}. Sin \texttt{return}, la función devuelve \texttt{None}.

\begin{formulabox}[Patrón clave]
\footnotesize
\texttt{resultado = funcion(argumento)} --- siempre guardar lo que la función devuelve en una variable, para poder usarlo después (imprimirlo, compararlo, sumarlo). Si la función usa \texttt{print()} en vez de \texttt{return}, esa variable queda en \texttt{None}.
\end{formulabox}

\begin{cajacodigo}
\begin{Verbatim}[fontsize=\footnotesize]
total = calcula_total(1500, 3)
print("Total:", total)
>> Total: 4500
\end{Verbatim}
\end{cajacodigo}

\textbf{Valores por omisión:} \texttt{def funcion(parametro=valor):} vuelve ese parámetro opcional al llamar --- siempre a la derecha de los parámetros obligatorios.

\begin{cajacodigo}
\begin{Verbatim}[fontsize=\footnotesize]
def calcula_total(precio, cantidad=1):
    return precio * cantidad

calcula_total(1500)      # cantidad = 1
calcula_total(1500, 3)   # cantidad = 3
\end{Verbatim}
\end{cajacodigo}

\end{multicols}

\end{document}
